% Generated draft: requires mathematical and source review.
\documentclass[11pt]{article}
\usepackage[margin=25mm]{geometry}
\usepackage{fontspec}
\setmainfont{Latin Modern Roman}
\usepackage{amsmath,amssymb,mathtools,unicode-math}
\setmathfont{Latin Modern Math}
\usepackage{enumitem,needspace,xurl,hyperref}
\setlength{\parindent}{0pt}
\setlength{\parskip}{5pt}
\setlength{\emergencystretch}{3em}
\newcommand{\N}{\mathbb{N}}
\newcommand{\Z}{\mathbb{Z}}
\newcommand{\Q}{\mathbb{Q}}
\newcommand{\R}{\mathbb{R}}
\newcommand{\C}{\mathbb{C}}
\newcommand{\F}{\mathbb{F}}
\newcommand{\A}{\mathbb{A}}
\newcommand{\PP}{\mathbb{P}}
\newcommand{\E}{\mathbb{E}}
\newcommand{\eps}{\varepsilon}
\newcommand{\dd}{\,\mathrm{d}}
\begin{document}
\section*{UnsolvedMath ID: 9400060}
\textbf{Source problem code:} catalogue entry
\subsection{Definitions and mathematical statement}
The source describes the Beilinson special-value conjectures for a motive, or a specified cohomological part of a smooth projective variety, at an appropriate integer argument. It does not supply the motive category, base field, weight, cohomological degree, integer range or regulator normalization, so the exact conjecture cannot be selected from the family.

Schematically, after those data are fixed, let \(L(M,s)\) be the intended motivic L-function and let \(r_{\mathrm{pred}}(M,m)\) be the rank predicted by the specified motivic-cohomology group at integer \(m\). The first assertion is
\[
 \operatorname{ord}_{s=m}L(M,s)=r_{\mathrm{pred}}(M,m).
\]
When the selected argument is a zero or regular point, define the leading coefficient by
\[
 L^*(M,m)=\lim_{s\to m}(s-m)^{-r_{\mathrm{pred}}(M,m)}L(M,s).
\]
The second assertion identifies this coefficient, up to the rational factor or rational-basis ambiguity prescribed by the original formulation, with the determinant of the appropriate Beilinson regulator. In standard geometric settings that regulator maps motivic cohomology to real Deligne-type cohomology. Its determinant is meaningful only after the relevant subspaces/quotients, dimensions, rational structures and period/trivialization factors have been specified. It is not the determinant of an unspecified map between arbitrary cohomology groups.

The excerpt does not identify which motivic group supplies the rank or whether analytic continuation, completed versus uncompleted L-functions, archimedean factors, or pole cases are included. These choices affect orders and coefficients. No unrestricted formula at every integer and no replacement of the prescribed rational normalization by an arbitrary numerical equality is inserted.

The supplied summary calls the general conjectures open and mentions Beilinson's \(K_3\)/elliptic-curve work, Deligne-type special values, and work of Flach, Burns, Nekov\'a\v{r} and others in elliptic-unit or Stark-type settings. It gives no precise hypotheses or citations, so these related results are not identified with a proof of every assertion in this family. The partial-progress status remains a source assessment pending primary-literature review.
\subsection{Short English statement}
At the intended integer, should a motivic L-function's vanishing order equal the specified motivic rank, and should its leading coefficient be the normalized regulator determinant? The source gives this broad family but omits the indices, argument range and normalization needed for an exact conjecture.
\subsection{Sources}
\begin{itemize}
\item Local registry: \url{data/unsolvedmath/problems.json}, permanent UnsolvedMath ID \(9400060\), source code catalogue entry; supplied statement, background, and research summary, updated 2026-08-06. Attributed to Alexander Beilinson.
\item Wikipedia contributors, \emph{List of unsolved problems in mathematics}, Number theory, analytic number theory bullet 60, pinned revision 1366636767: \url{https://en.wikipedia.org/w/index.php?title=List_of_unsolved_problems_in_mathematics&oldid=1366636767#Number_theory}. Supplied MediaWiki-source extraction plus linked special-values-page resolution, accessed 2026-07-29.
\item The summary refers to the Wikipedia ``Beilinson conjectures'' page and names Beilinson, Deligne, Flach, Burns, Nekov\'a\v{r} and related special cases without full original-paper citations or exact statements.
\end{itemize}
\end{document}
